\part{Anyons, Topology, and Generalized Theories}

\chapter{Planar exchange produces braid groups}
\label{ch:braids}

\physicalquestion{Why can exchanging identical particles in two dimensions
retain information that disappears in three dimensions?}

\modelassumptions{Particles are pointlike, identical, and constrained to a
two-dimensional region.  Collision configurations are excluded, motion is
adiabatic, and worldlines are compared up to deformation that fixes their
endpoints.}

\section{Which group records planar worldlines?}

For $n$ ordered particles in a surface $\Sigma$, remove collisions from
$\Sigma^n$ and quotient by permutations to obtain the configuration space of
identical particles.  Its fundamental group is the braid group $B_n$ when
$\Sigma$ is a disk.  It has generators $\sigma_1,\ldots,\sigma_{n-1}$ and
relations
\[
\sigma_i\sigma_{i+1}\sigma_i
=\sigma_{i+1}\sigma_i\sigma_{i+1},
\qquad
\sigma_i\sigma_j=\sigma_j\sigma_i\quad(|i-j|\ge2).
\]
Unlike the symmetric group, $B_n$ does not impose $\sigma_i^2=1$.

A monoidal category has a tensor product, a unit, and coherent associativity
maps.  A braiding is a natural family
$c_{x,y}:x\otimes y\to y\otimes x$ satisfying the two hexagon identities.

\section{Why does categorical braiding act on many-particle states?}

\begin{theorem}[Braiding represents the braid group]
For any object $x$ in a braided monoidal category, the maps
\[
\rho(\sigma_i)
=1_x^{\otimes(i-1)}\otimes c_{x,x}\otimes
1_x^{\otimes(n-i-1)}
\]
define a representation of $B_n$ on $x^{\otimes n}$, with associators
inserted coherently.
\end{theorem}

\begin{proof}
Naturality and the hexagon identities give the Yang--Baxter equality on three
adjacent tensor factors, which is the first braid relation.  Maps acting on
disjoint tensor factors commute, giving the second.  Mac Lane coherence makes
the result independent of how the iterated tensor product is parenthesized.
\end{proof}

\section{What physical information survives an exchange?}

Adiabatic transport around a braid may act on the degenerate state space by
$\rho(\beta)$.  In an abelian anyon sector this action is a phase; in a
nonabelian sector it can be a noncommuting matrix.  Three-dimensional exchange
adds a homotopy that collapses the braid to a permutation, forcing the more
familiar boson/fermion alternatives under standard assumptions.  Planar
topology removes that collapse and therefore demands braid data.

\section{What further contact should be proved?}

\begin{exercise}[The symmetric group is a quotient]
Prove that imposing $\sigma_i^2=1$ on $B_n$ gives the symmetric group $S_n$.
Interpret the added relation as forgetting overcrossing versus undercrossing.
\end{exercise}

\begin{exercise}[Temperley--Lieb generators braid]
Given the relations in Chapter~\ref{ch:subfactors}, determine scalars $a,b$
for which $a1+be_i$ satisfy the braid relations.
\end{exercise}


\chapter{Fusion requires tensor categories}
\label{ch:fusion}

\physicalquestion{When anyons are combined, why can the result have several
possible particle types and an internal fusion space?}

\modelassumptions{There are finitely many topological charge types, fusion
spaces are finite dimensional, every charge has an antiparticle, and local
microscopic excitations have been quotiented out.}

\section{Which language contains fusion and exchange together?}

In a semisimple rigid $\C$-linear monoidal category, simple objects represent
anyon types.  Fusion is the decomposition
\[
x\otimes y\cong\bigoplus_z N_{xy}^z z,
\qquad
\Hom(z,x\otimes y)
\]
is the corresponding fusion space.  Rigidity supplies an antiparticle $x^*$.
An associator compares the two orders of fusing three particles, while a
braiding exchanges them.  Pentagon and hexagon coherence express consistency
of all composite processes.

The monodromy $M_{x,y}=c_{y,x}c_{x,y}$ is a full winding of $x$ around $y$.

\section{How is abelian monodromy classified?}

\begin{theorem}[Pointed monodromy]
Suppose every simple object is invertible and the simple classes form the
cyclic group $\Z/m$ generated by $a$.  A scalar monodromy compatible with
fusion has the form
\[
M_{a^r,a^s}=\exp\!\left(\frac{2\pi i krs}{m}\right)
\]
for some $k\in\Z/m$.
\end{theorem}

\begin{proof}
The hexagon identities make monodromy multiplicative in either variable, so
$M$ is a bicharacter on $\Z/m\times\Z/m$.  It is determined by
$\zeta=M_{a,a}$.  Since $a^m\cong1$, multiplicativity gives $\zeta^m=1$;
write $\zeta=e^{2\pi ik/m}$.  Bilinearity then gives the formula.
\end{proof}

\section{What phase results for one-third anyons?}

For $m=3$ and $k=1$, winding $a$ once around $a$ multiplies the state by
$e^{2\pi i/3}$.  In an interferometer, changing the number of enclosed anyons
therefore advances the interference phase by $2\pi/3$ in the ideal abelian
model.  Real devices also include Coulomb coupling, edge reconstruction, and
decoherence; the categorical prediction concerns the topological contribution
after those effects are controlled \cite{nakamura}.

\section{What further contact should be proved?}

\begin{exercise}[A three-cocycle controls reassociation]
For $G$-graded vector spaces, multiply the associator on homogeneous degrees
$g,h,k$ by $\omega(g,h,k)\in U(1)$.  Prove that the pentagon is exactly the
$3$-cocycle equation and that changing $\omega$ by a coboundary gives a
monoidally equivalent category.
\end{exercise}

\begin{exercise}[Fusion dimensions]
If positive numbers $d_x$ satisfy
$d_xd_y=\sum_zN_{xy}^zd_z$, prove that $d$ is a positive character of the
fusion ring.  Compute it for the Fibonacci rule $\tau\otimes\tau=1\oplus\tau$.
\end{exercise}


\chapter{Chern--Simons theory turns braids into amplitudes}
\label{ch:chern-simons}

\physicalquestion{Which field theory assigns an amplitude to a braid using
only topology rather than a spacetime metric?}

\modelassumptions{Spacetime is a closed oriented three-manifold, the gauge
group is compact, and the level is chosen so the exponentiated action is gauge
invariant.  Path-integral statements are interpreted formally unless a
rigorous TQFT construction is specified.}

\section{Which action sees knots but not distances?}

For a connection $A$ on a principal $G$-bundle, the Chern--Simons action is
\[
S_{\mathrm{CS}}(A)
=\frac{k}{4\pi}\int_M
\tr\!\left(A\wedge dA+\frac23A\wedge A\wedge A\right).
\]
It uses orientation and a trace pairing but no metric.  Variation gives
$F_A=0$, so classical solutions are flat connections whose holonomy may still
be globally nontrivial.

A labeled closed worldline $K$ defines a Wilson loop
$W_R(K)=\tr_R\operatorname{Hol}_A(K)$.

\section{Why must the level be quantized?}

\begin{theorem}[Large-gauge invariance]
With the standard normalization for compact simple simply connected $G$, a
gauge transformation $g:M\to G$ changes the Chern--Simons action by
$2\pi k\,\deg(g)$.  Hence $e^{iS_{\mathrm{CS}}}$ is gauge invariant for all
gauge transformations exactly when $k\in\Z$.
\end{theorem}

\begin{proof}
Substitute $A^g=g^{-1}Ag+g^{-1}dg$ into the action.  Cyclicity of trace
separates an exact transgression term and
\[
\frac{k}{12\pi}\int_M\tr(g^{-1}dg)^3.
\]
The exact term integrates to zero on closed $M$.  Under the standard trace
normalization, the remaining integral is $2\pi k$ times the integer winding
number of $g$.  The exponential is therefore invariant for all winding
numbers precisely at integral level.
\end{proof}

\section{How do Wilson loops recover anyon data?}

In Chern--Simons TQFT, exchanging worldlines evaluates the braid morphism of a
related modular tensor category, and closing a braid evaluates its
categorical trace.  For suitable $SU(2)$ level, Wilson-loop amplitudes produce
Jones-type link invariants.  Thus fusion and braiding are not appended to the
field theory: they are the invariant data assigned to punctures and
worldlines.  The formal path integral motivates this relation; rigorous
Reshetikhin--Turaev constructions implement it categorically.
\cite{witten-jones,reshetikhin-turaev}

\section{What further contact should be proved?}

\begin{exercise}[Flat does not mean trivial]
Construct flat $U(1)$ connections with distinct holonomy on a three-torus and
compare with Chapter~\ref{ch:ab}.
\end{exercise}

\begin{exercise}[Wilson loops respect braid closure]
Assuming the ribbon-category evaluation rules, prove that isotopic closed
braids have equal amplitudes and identify where framing enters.
\end{exercise}


\chapter{Functorial theories stabilize observables}
\label{ch:generalized}

\physicalquestion{What common structure turns fusion, gluing, bundles, and
Fredholm operators into deformation-stable numerical observables?}

\modelassumptions{Categories are additive where Grothendieck completion is
used, tensor products distribute over direct sums, and generalized cohomology
is evaluated on finite CW complexes unless stated otherwise.}

\section{How are extended processes composed?}

A strong monoidal functor $F:\mathcal C\to\mathcal D$ comes with coherent
isomorphisms $F(x)\otimes F(y)\cong F(x\otimes y)$ and preserves the tensor
unit.  A three-dimensional TQFT is a symmetric monoidal functor from a bordism
category to vector spaces: disjoint union becomes tensor product and gluing
becomes composition.

For an additive monoidal category, its Grothendieck group is generated by
objects with $[x\oplus y]=[x]+[y]$ and multiplication
$[x][y]=[x\otimes y]$.  This is decategorification: fusion spaces and
morphisms are forgotten, while stable addition and multiplication remain.

\section{What survives a monoidal passage?}

\begin{theorem}[Monoidal decategorification]
An additive strong monoidal functor induces a ring homomorphism
\[
K_0(F):K_0(\mathcal C)\longrightarrow K_0(\mathcal D).
\]
If the categories are rigid and traces are preserved, categorical dimensions
and traces of closed processes are preserved as well.
\end{theorem}

\begin{proof}
Additivity gives
$[F(x\oplus y)]=[F(x)]+[F(y)]$, so the assignment $[x]\mapsto[F(x)]$
respects the defining Grothendieck relation.  The monoidal comparison gives
$[F(x\otimes y)]=[F(x)\otimes F(y)]$, hence multiplication is preserved.
Strong monoidal functors carry evaluation and coevaluation maps to duality
maps, so the composite defining categorical trace maps to the corresponding
composite in $\mathcal D$.
\end{proof}

\section{How do generalized theories unify the earlier invariants?}

A reduced generalized cohomology theory is homotopy invariant, sends cofiber
sequences to long exact sequences, and has suspension isomorphisms.  Complex
K-theory is two-periodic and admits equivalent bundle, projection, unitary,
and Fredholm models.  Those models connect the occupied Bloch bundle of
Chapter~\ref{ch:hall} to the index of Chapter~\ref{ch:index}.

The recurring pattern is now explicit:
\[
\begin{gathered}
\text{physical structure}\longrightarrow\text{functorial stable class}
\\[-2pt]
\longrightarrow\text{pairing or trace}\longrightarrow\text{observable}.
\end{gathered}
\]
Aharonov--Bohm phase, Hall conductance, Fredholm index, and anyon amplitude
use different categories, but each is protected because the middle class is
unchanged by the allowed deformation.

\section{What further contact should be proved?}

\begin{exercise}[Cohomology preserves supertraces]
Prove that cohomology is symmetric monoidal on bounded complexes of
finite-dimensional complex vector spaces and deduce the Hopf trace identity.
\end{exercise}

\begin{exercise}[Part VI synthesis]
Start with the pointed $\Z/3$ braided category of
Chapter~\ref{ch:fusion}.  Construct its braid representations, close a braid
using categorical trace, interpret the result as a Wilson-loop amplitude, and
decategorify its fusion rules.  Identify which data are lost at each passage
and verify that the $2\pi/3$ monodromy phase remains.
\end{exercise}

\begin{exercise}[Final prerequisite audit]
For one invariant from each part, list every mathematical structure on which
its definition depends.  Order those prerequisites independently of the
book's narrative and verify that every dependency appears before its first
use.  If an assumption belongs to a physical model rather than a theorem,
mark it explicitly.
\end{exercise}

\begin{quote}\textbf{Closing perspective.} Superposition demanded linear
spaces; invariant transformations demanded representations; local comparison
demanded geometry; robust global response demanded topology; matrix-level
composition demanded noncommutative analysis; and fusion with exchange
demanded categories.  At each stage, the failure of the preceding framework
selected the next coordinate-free structure.\end{quote}
